\documentclass{article}
\usepackage{../assignment_preamble}

\title{Homework 7}
\author{Ravi Kini}
\date{November 26, 2025}

\begin{document}

\maketitle

\problem{}
We test the counterfactual probability $\theta' = \frac{1}{3}$ against the true probability of $\theta = \frac{1}{2}$. The learning algorithm learns the true probability when the observed frequency is within $\frac{|1/2 - 1/3|}{2} = \frac{1}{12}$ of $\theta$. Under the IID assumption, with $n = 150$ samples, this occurs with probability at least:
\begin{align*}
    P_{1/2, 150}(|\text{observed frequency} - 1/2| \leq \frac{1}{12}) \geq 1 - \frac{1}{4(150){(1/12)}^2} = 0.76
\end{align*}
Similarly, with $n = 500$ samples:
\begin{align*}
    P_{1/2, 500}(|\text{observed frequency} - 1/2| \leq \frac{1}{12}) \geq 1 - \frac{1}{4(500){(1/12)}^2} = 0.928
\end{align*}

\clearpage
\problem{}
Consider the Bayesian network with the following DAG over Boolean variables:
\begin{center}
\begin{tikzpicture}[
       decoration = {markings,
                     mark=at position .5 with {\arrow{Stealth[length=2mm]}}},
       dot/.style = {circle, fill, inner sep=2.4pt, node contents={},
                     label=#1},
every edge/.style = {draw, postaction=decorate}
                        ]
\node (x) at (0,0) [dot=above:$X$];
\node (u) at (-2,-1) [dot=left:$U$];
\node (y) at (0,-2) [dot=right:$Y$];
%
\path
(x) edge (u)    (u) edge (y)
(x) edge (y);
\end{tikzpicture}
\end{center}
with the following probability tables:
\begin{center}
\begin{tabular}{|c|c|}
    \hline
    $P(X = 0)$ & $P(X = 1)$ \\
    \hline
    $1/2$ & $1/2$ \\
    \hline
\end{tabular}
\end{center}
\begin{center}
\begin{tabular}{|c|c|c|}
    \hline
    $x$ & $P(U = 0 | X = x)$ & $P(U = 1 | X = x)$ \\
    \hline
    $0$ & $9/10$ & $1/10$ \\
    \hline
    $1$ & $1/5$ & $4/5$ \\
    \hline
\end{tabular}
\end{center}
\begin{center}
\begin{tabular}{|c|c|c|c|}
    \hline
    $u$ & $x$ & $P(Y = 0 | U = u, X = x)$ & $P(Y = 1 | U = u, X = x)$ \\
    \hline
    $0$ & $0$ & $1$ & $0$ \\
    \hline
    $0$ & $1$ & $1$ & $0$ \\
    \hline
    $1$ & $0$ & $1 - q$ & $q$ \\
    \hline
    $1$ & $1$ & $1 - (1 - r)q$ & $(1 - r)q$ \\
    \hline
\end{tabular}
\end{center}
We first compute $P(Y = y|X = x)$.
\begin{align*}
    P(Y = y|X = x) & = P(Y = y, U = 0|X = x) + P(Y = y, U = 1|X = x) \\
    & = P(Y = y|U = 0, X = x)P(U = 0, X = x) \\
    & + P(Y = y|U = 1, X = x)P(U = 1, X = x)
\end{align*}
\begin{center}
\begin{tabular}{|c|c|c|}
    \hline
    $x$ & $P(Y = 0 | X = x)$ & $P(Y = 1 | X = x)$ \\
    \hline
    $0$ & $1 - q/10$ & $q/10$ \\
    \hline
    $1$ & $1 - 4(1 - r)q/5$ & $4(1 - r)q/5$ \\
    \hline
\end{tabular}
\end{center}
Similarly, we compute $P(Y)$.
\begin{align*}
    P(Y = y) & = P(Y = y, X = 0) + P(Y = y, X = 1) \\
    & = P(Y = y | X = 0)P(X = 0) + P(Y = y | X = 1)P(X = 1)
\end{align*}
\begin{center}
\begin{tabular}{|c|c|}
    \hline
    $P(X = 0)$ & $P(X = 1)$ \\
    \hline
    $1 - q/20 - 2(1 - r)q/5$ & $q/20 + 2(1 - r)q/5$ \\
    \hline
\end{tabular}
\end{center}
Setting $P(Y = 1|X = 0)$ and $P(Y = 1)$ equal:
\begin{align*}
    q/10 & = q/20 + 2(1 - r)q/5 \\
    q/20 - 2(1 - r)q/5 & = 0 \\
    q(1/20 - 2(1 - r)/5) & = 0 \\
    q(r - 7/8) & = 0 \\
\end{align*}
Either $q = 0$ (and $r$ takes on any value) or $r = 7/8$ (and $q$ takes on any value).

\clearpage
\problem{}
Consider the Bayesian network $N_1$ with the following DAG over Boolean variables:
\begin{center}
\begin{tikzpicture}[
       decoration = {markings,
                     mark=at position .5 with {\arrow{Stealth[length=2mm]}}},
       dot/.style = {circle, fill, inner sep=2.4pt, node contents={},
                     label=#1},
every edge/.style = {draw, postaction=decorate}
                        ]
\node (x) at (0,0) [dot=above:$X$];
\node (z) at (0,-1) [dot=left:$Z$];
\node (y) at (0,-2) [dot=right:$Y$];
%
\path
(x) edge (z)    (z) edge (y);
\end{tikzpicture}
\end{center}
with the following probability tables:
\begin{center}
\begin{tabular}{|c|c|}
    \hline
    $P(X = 0)$ & $P(X = 1)$ \\
    \hline
    $1/2$ & $1/2$ \\
    \hline
\end{tabular}
\end{center}
\begin{center}
\begin{tabular}{|c|c|c|}
    \hline
    $x$ & $P(Z = 0 | X = x)$ & $P(Z = 1 | X = x)$ \\
    \hline
    $0$ & $1/3$ & $2/3$ \\
    \hline
    $1$ & $2/3$ & $1/3$ \\
    \hline
\end{tabular}
\end{center}
\begin{center}
\begin{tabular}{|c|c|c|}
    \hline
    $z$ & $P(Y = 0 | Z = z)$ & $P(Y = 1 | Z = z)$ \\
    \hline
    $0$ & $1/6$ & $5/6$ \\
    \hline
    $1$ & $5/6$ & $1/6$ \\
    \hline
\end{tabular}
\end{center}
We then compute $P(Y = y, Z = z, X = x)$.
\begin{align*}
    P(Y = y, Z = z, X = x) & = P(Y = y|Z = z, X = x)P(Z = z, X = x) \\
    & = P(Y = y|Z = z)P(Z = z, X = x) \\
    & = P(Y = y|Z = z)P(Z = z | X = x)P(X = x)
\end{align*}
Constructing the joint probability distribution:
\begin{center}
\begin{tabular}{|c|c|c|c|}
    \hline
    $y$ & $z$ & $x$ & $P(Y = y, Z = z, X = x)$ \\
    \hline
    $0$ & $0$ & $0$ & $1/6 \cdot 1/3 \cdot 1/2 = 1/36$ \\
    \hline
    $0$ & $0$ & $1$ & $1/6 \cdot 2/3 \cdot 1/2 = 1/18$ \\
    \hline
    $0$ & $1$ & $0$ & $5/6 \cdot 2/3 \cdot 1/2 = 5/18$ \\
    \hline
    $0$ & $1$ & $1$ & $5/6 \cdot 1/3 \cdot 1/2 = 5/36$ \\
    \hline
    $1$ & $0$ & $0$ & $5/6 \cdot 1/3 \cdot 1/2 = 5/36$ \\
    \hline
    $1$ & $0$ & $1$ & $5/6 \cdot 2/3 \cdot 1/2 = 5/18$ \\
    \hline
    $1$ & $1$ & $0$ & $1/6 \cdot 2/3 \cdot 1/2 = 1/18$ \\
    \hline
    $1$ & $1$ & $1$ & $1/6 \cdot 2/3 \cdot 1/2 = 1/36$ \\
    \hline
\end{tabular}
\end{center}

\clearpage
\problem{}
Consider the Bayesian network $N_2$ with the following DAG over Boolean variables:
\begin{center}
\begin{tikzpicture}[
       decoration = {markings,
                     mark=at position .5 with {\arrow{Stealth[length=2mm]}}},
       dot/.style = {circle, fill, inner sep=2.4pt, node contents={},
                     label=#1},
every edge/.style = {draw, postaction=decorate}
                        ]
\node (x) at (0,0) [dot=above:$X$];
\node (y) at (-2,-1) [dot=left:$Y$];
\node (z) at (0,-2) [dot=right:$Z$];
%
\path
(x) edge (y)    (y) edge (z)
(x) edge (z);
\end{tikzpicture}
\end{center}
Note that:
\begin{align*}
    P(Y = y | X = x) & = \frac{P(Y = y, X = x)}{P(X = x)} \\
    P(Z = z | Y = y, X = x) & = \frac{P(Z = z, Y = y, X = x)}{P(Y = y, X = x)}
\end{align*}
We can then construct the following probability tables:
\begin{center}
\begin{tabular}{|c|c|}
    \hline
    $P(X = 0)$ & $P(X = 1)$ \\
    \hline
    $1/2$ & $1/2$ \\
    \hline
\end{tabular}
\end{center}
\begin{center}
\begin{tabular}{|c|c|c|}
    \hline
    $x$ & $P(Y = 0 | X = x)$ & $P(Y = 1 | X = x)$ \\
    \hline
    $0$ & $11/18$ & $7/18$ \\
    \hline
    $1$ & $7/18$ & $11/18$ \\
    \hline
\end{tabular}
\end{center}
\begin{center}
\begin{tabular}{|c|c|c|c|}
    \hline
    $y$ & $x$ & $P(Z = 0 | Y = y, X = x)$ & $P(Z = 1 | Y = y, X = x)$ \\
    \hline
    $0$ & $0$ & $1/11$ & $10/11$ \\
    \hline
    $0$ & $1$ & $2/7$ & $5/7$ \\
    \hline
    $1$ & $0$ & $5/7$ & $2/7$ \\
    \hline
    $1$ & $1$ & $10/11$ & $1/11$ \\
    \hline
\end{tabular}
\end{center}
We then compute $P(Z = z, Y = y, X = x)$:
\begin{align*}
    P(Z = z, Y = y, X = x) & = P(Z = z|Y = y, X = x)P(Y = y, X = x) \\
    & = P(Z = z|Y = y, X = x)P(Y = y | X = x)P(X = x) \\
\end{align*}
Constructing the joint probability distribution:
\begin{center}
\begin{tabular}{|c|c|c|c|}
    \hline
    $z$ & $y$ & $x$ & $P(Z = z, Y = y, X = x)$ \\
    \hline
    $0$ & $0$ & $0$ & $1/11 \cdot 11/18 \cdot 1/2 = 1/36$ \\
    \hline
    $0$ & $0$ & $1$ & $2/7 \cdot 7/18 \cdot 1/2 = 1/18$ \\
    \hline
    $0$ & $1$ & $0$ & $5/7 \cdot 7/18 \cdot 1/2 = 5/36$ \\
    \hline
    $0$ & $1$ & $1$ & $10/11 \cdot 11/18 \cdot 1/2 = 5/18$ \\
    \hline
    $1$ & $0$ & $0$ & $10/11 \cdot 11/18 \cdot 1/2 = 5/18$ \\
    \hline
    $1$ & $0$ & $1$ & $5/7 \cdot 7/18 \cdot 1/2 = 5/36$ \\
    \hline
    $1$ & $1$ & $0$ & $2/7 \cdot 7/18 \cdot 1/2 = 1/18$ \\
    \hline
    $1$ & $1$ & $1$ & $1/11 \cdot 11/18 \cdot 1/2 = 1/36$ \\
    \hline
\end{tabular}
\end{center}

\clearpage
\problem{}
Note that $X$ and $Y$ are independent in $N_1$. Since $N_1$ and $N_2$ have the same joint proability distribution table, this also holds in $N_2$. However, there is an edge between $X$ and $Y$ in $N_2$, so this does not hold for all Bayesian networks with the same DAG. Consequently, $N_2$ is not faithful.
\end{document}